% Generated draft: requires mathematical and source review.
\documentclass[11pt]{article}
\usepackage[margin=25mm]{geometry}
\usepackage{fontspec}
\setmainfont{Latin Modern Roman}
\usepackage{amsmath,amssymb,mathtools,unicode-math}
\setmathfont{Latin Modern Math}
\usepackage{enumitem,needspace,xurl,hyperref}
\setlength{\parindent}{0pt}
\setlength{\parskip}{5pt}
\setlength{\emergencystretch}{3em}
\newcommand{\N}{\mathbb{N}}
\newcommand{\Z}{\mathbb{Z}}
\newcommand{\Q}{\mathbb{Q}}
\newcommand{\R}{\mathbb{R}}
\newcommand{\C}{\mathbb{C}}
\newcommand{\F}{\mathbb{F}}
\newcommand{\A}{\mathbb{A}}
\newcommand{\PP}{\mathbb{P}}
\newcommand{\E}{\mathbb{E}}
\newcommand{\eps}{\varepsilon}
\newcommand{\dd}{\,\mathrm{d}}
\begin{document}
\section*{UnsolvedMath ID: 701463}
\textbf{Source problem code:} catalogue entry
\subsection{Definitions and mathematical statement}
A pro-$p$ group, for a prime $p$, is an inverse limit of finite $p$-groups with their discrete topologies. A pronilpotent group is an inverse limit of finite nilpotent groups with their discrete topologies; these inverse limits carry their profinite topologies.

For a topological group $G$, consider a finite subnormal series of closed subgroups
\[
S:\quad G_0\trianglelefteq G_1\trianglelefteq\cdots\trianglelefteq G_n=G.
\]
Thus each $G_i$ is closed in $G$, and $G_{i-1}$ is normal in $G_i$ for $1\leq i\leq n$. Define the infinite-length of $S$ by
\[
\ell_\infty(S)=\bigl|\{i\in\{1,\ldots,n\}:G_i/G_{i-1}\text{ is infinite}\}\bigr|,
\]
and define the virtual length of $G$ by
\[
\operatorname{vl}(G)=\sup_S\ell_\infty(S)\in\Z_{\geq0}\cup\{\infty\},
\]
where $S$ ranges over all such finite series. In particular, finite virtual length means that there is a single integer $M\geq0$ bounding the number of infinite factors in every such series. The supplied definition does not explicitly require $G_0=\{1\}$; the definition above follows its displayed formulation, and the intended initial-term convention remains a point for source review.

A closed subgroup $H\leq G$ is subnormal in $G$ if there is a finite chain of closed subgroups
\[
H=H_0\trianglelefteq H_1\trianglelefteq\cdots\trianglelefteq H_r=G,
\]
where $r\geq0$. It is topologically finitely generated if there is a finite set $F\subseteq H$ such that $H=\overline{\langle F\rangle}$, with closure taken in $H$. For a closed subgroup, requiring closed intermediate terms agrees with the usual finite-chain notion of subnormality: intermediate terms can be replaced by their closures.

\textbf{Problem.} For every prime $p$ and every pro-$p$ group $G$, is the following implication true?
\[
\left[
\begin{gathered}
\text{Every closed subnormal subgroup }H\leq G\\
\text{is topologically finitely generated}
\end{gathered}
\right]
\quad\Longrightarrow\quad
\operatorname{vl}(G)<\infty.
\]
More generally, does this implication hold for every pronilpotent group $G$?

\textbf{Supplied background.} The record attributes to N. Gavioli, V. Monti, and C. M. Scoppola the converse implication for pronilpotent groups: finite virtual length implies that every closed subnormal subgroup is topologically finitely generated. This is supplied background attribution, not an independently verified theorem or literature assessment.
\subsection{Short English statement}
The virtual length counts, with a uniform bound over finite closed subnormal series, how many quotient factors can be infinite. The question asks whether topological finite generation of every closed subnormal subgroup forces such a bound for every pro-$p$ group, and more generally for every pronilpotent group. The supplied registry labels the problem open; that status has not been independently verified.
\subsection{Sources}
\begin{itemize}
\item Local UnsolvedMath registry, permanent ID \texttt{701463}, exact source problem code \texttt{catalogue entry}, title \emph{Kourovka Notebook Problem 18.34}. The supplied record labels the problem \emph{open}; its update date is \texttt{2026-10-16}. This is supplied registry metadata, not independent source verification.
\item \emph{The Kourovka Notebook (Unsolved Problems in Group Theory)}, original problem 18.34, 18th issue (2014); also identified in the record with the 21st issue edition (2026), \texttt{arXiv:1401.0300v48}. Proposed by N. Gavioli, V. Monti, and C. M. Scoppola. The supplied collection description names E. I. Khukhro and V. D. Mazurov as editors.
\item N. Gavioli, V. Monti, and C. M. Scoppola, \emph{J. Austral. Math. Soc.} \textbf{95}, no. 3 (2013), 343--355. The record cites this reference for the background implication concerning pronilpotent groups of finite virtual length. No article title, DOI, or URL is supplied.
\end{itemize}
\end{document}
